% Обзор литературы 0002.
% Сборка: tectonic review.tex (или xelatex + bibtex)
% PDF кладётся в out/, в репозиторий не коммитится.
\documentclass[11pt,a4paper]{article}

\usepackage{iftex}
\ifPDFTeX
  \usepackage[T2A]{fontenc}
  \usepackage[utf8]{inputenc}
  \usepackage[russian]{babel}
\else
  \usepackage{fontspec}
  \usepackage[russian]{babel}
  \IfFontExistsTF{Times New Roman}
    {\setmainfont{Times New Roman}}
    {\IfFontExistsTF{DejaVu Serif}
      {\setmainfont{DejaVu Serif}}
      {\IfFontExistsTF{Liberation Serif}
        {\setmainfont{Liberation Serif}}
        {}}}
\fi

\usepackage[margin=2.4cm]{geometry}
\usepackage{amsmath}
\usepackage{amssymb}
\usepackage{booktabs}
\usepackage{longtable}
\usepackage{array}
\usepackage{enumitem}
\usepackage{hyperref}
\usepackage{microtype}
\setlength{\emergencystretch}{4em}

\hypersetup{
  colorlinks=true,
  linkcolor=black,
  citecolor=black,
  urlcolor=blue
}

\setlist[itemize]{leftmargin=1.4em,itemsep=0.2em,topsep=0.2em}
\setlist[enumerate]{leftmargin=1.6em,itemsep=0.2em,topsep=0.2em}

\title{Обнаружение и коррекция неверной оценки сложности\\
по ходу выполнения задачи: обзор литературы}
\author{Обзор 0002}
\date{8 ноября 2026}

\begin{document}
\maketitle

\begin{abstract}
\noindent
Обзор отвечает на вопрос, как во время исполнения задачи понять, что её
сложность была оценена неверно, и как скорректировать оценку и план.
Включено 50 работ: онлайн-мониторинг LLM-агентов, калибровка уверенности,
обнаружение аномалий и зависаний, перепланирование, бюджетные политики,
пределы самокоррекции, а также классические аналоги — мониторинг бизнес-
процессов, earned value management и обнаружение дрейфа. Вывод: сигналов
расхождения много, и они работают; гораздо слабее изучено решение о
коррекции. Точность монитора сама по себе не говорит, стоит ли вмешиваться:
ценность определяется ожидаемой выгодой от действия. Надёжная схема состоит
из трёх слоёв: последовательный тест по остаткам прогноза с явным
компромиссом «ложная тревога — задержка», выбор действия по решающей
ценности и обновление оценщика по итогам прогона. Полагаться на вербальную
самокоррекцию модели нельзя; нужен внешний сигнал и явное правило. Готовой
интегрированной схемы для кодового агента при общем бюджете токенов
литература не даёт.
\end{abstract}

\section{Вопрос и зачем}

\subsection{Вопрос}

Как во время исполнения задачи понять, что её сложность была оценена
неверно, и как скорректировать оценку и план, не потратив остаток бюджета
впустую?

Вопрос распадается на три:

\begin{enumerate}
  \item \textbf{Что наблюдать.} Какой сигнал по ходу исполнения говорит о
        том, что прогноз стоимости, успеха или выполнимости расходится с
        фактом?
  \item \textbf{Как решить.} Когда сигнала достаточно, чтобы объявить
        оценку неверной, с учётом ложных тревог и задержки обнаружения?
  \item \textbf{Что исправлять.} Обновить число, перепланировать, сменить
        стратегию, остановиться или эскалировать — и как это сделать в
        пределах остатка бюджета?
\end{enumerate}

Слово «сложность» здесь, как и в обзоре 0001, означает не одно число, а
набор прогнозов: вероятность успеха $p$, распределение стоимости $C$,
остаток работы $R$ и выполнимость $F$. Ошибка оценивания — это
расхождение любого из них с наблюдением, достаточно большое, чтобы
изменить оптимальное действие.

\subsection{Зачем}

Обзор 0001 показал, как оценить задачу до исполнения, и одновременно
почему одной априорной оценки мало. TokenCast~\cite{ouyang2026tokencast}
сообщает, что расход токенов на одну и ту же задачу различается более чем
на порядок между прогонами. BAGEN~\cite{lin2026bagen} находит, что
сильные агенты систематически оптимистичны и продолжают тратить бюджет на
задачах, которые уже вряд ли будут решены. Значит, оценка должна
пересматриваться по ходу дела.

Для харнесса это практический вопрос. Сейчас харнесс не оценивает задачи
заранее и не пересматривает решение: задача решается до конца либо бюджет
кончается. Если добавить оценку, но не добавить коррекцию, харнесс будет
уверенно идти по неверному плану. Ответ на вопрос определяет, какие
сигналы, пороги и действия закладывать в изменение
\texttt{04\_budget\_and\_selection.md}, и стоит ли вообще писать это
изменение до эксперимента.

\subsection{Термины и границы}

\textbf{Ошибка оценки} (estimation error) — прогноз не совпал с фактом.
\textbf{Ошибка исполнения} (execution error) — план был верным, но агент
ошибся в шаге. \textbf{Невыполнимость} — задача не решается данным агентом
при любом разумном бюджете. Это разные причины, и коррекция у них разная:
в первом случае обновляют прогноз, во втором перепланируют, в третьем
останавливаются или эскалируют. Литература часто смешивает их, потому что
все три проявляются как «что-то пошло не так».

\textbf{Коррекция} рассматривается на двух горизонтах:

\begin{itemize}
  \item \emph{внутри прогона} — обновить прогноз и сменить действие на
        основании частичной траектории;
  \item \emph{между прогонами} — переобучить или перекалибровать оценщик
        по завершённым траекториям, чтобы следующая априорная оценка была
        точнее.
\end{itemize}

Обзор не проектирует изменение. Он фиксирует, какие сигналы проверены,
какие правила коррекции измерялись и где остаётся пробел.

\section{Метод}

\subsection{Базы и запросы}

Поиск выполнен 8 ноября 2026 года в arXiv API (поле \texttt{abs:}) и
OpenAlex (фильтр по заголовку и аннотации). Просматривались
верхние 1--30 записей на запрос, сортировка по релевантности. Язык —
английский, годы — 2004--2026; более ранние работы включались только как
метод, на который прямо опираются современные. Точные запросы, даты и
числа — в \texttt{protocol.md}.

Semantic Scholar и Google Scholar не использовались: в прошлом обзоре они
вернули HTTP~403. Это ограничивает полноту, особенно индустриальными
отчётами.

\subsection{Критерии отбора}

Включались работы, которые отвечают хотя бы на один из трёх вопросов:
\begin{itemize}
  \item дают сигнал расхождения прогноза и факта, доступный во время
        исполнения;
  \item дают правило выбора действия после обнаружения расхождения;
  \item дают механизм обновления оценки или политики оценивания.
\end{itemize}

Допускались LLM-агенты, кодовые агенты, планирование, мониторинг бизнес-
процессов и управление проектами. Классические работы включались, если их
сигнал или правило переносится на агента. Исключались работы, где сложность
измеряется только постфактум, где нужны недоступные сигналы (сенсоры,
скрытые состояния другой модели) и где нет метода или количественной
оценки.

\subsection{Отбор}

Идентифицировано 880 записей, после удаления 10 дубликатов просмотрено 870.
На скрининге исключено 792 записи, для детального разбора отобрано 78,
после разбора исключено 28. В обзор включено 50 работ. Схема и причины
исключения — в \texttt{sources.md}.

\subsection{Извлечение и оценка качества}

Для каждой работы выписаны: что обнаруживается, сигналы, метод, данные и
результат, ограничение. Надёжность оценивалась по рецензируемости, размеру
данных, наличию отложенной выборки или независимой сертификации, сравнению
с baseline, явной обработке ложных тревог и задержки. Препринты и
Zenodo-отчёты помечены как менее надёжные. Мета-анализ не проводился:
метрики несопоставимы.

\section{Что известно}

\subsection{Ошибка оценки видна как расхождение прогноза и факта}

Самое прямое определение ошибки: наблюдение вышло за пределы прогноза.
Работы этой группы учат прогнозировать стоимость и обновлять прогноз по
мере исполнения.

\textbf{Стоимость надо прогнозировать распределением, а не числом.}
TokenCast~\cite{ouyang2026tokencast} представляет каждый сегмент
траектории парой «собственный расход» и «прирост контекста», а затем
композирует сегменты в кумулятивную оценку. Новые наблюдения обновляют
прогноз без дополнительных вызовов LLM; среднее время обновления —
32{,}8~мс на прогон SWE-bench Verified. На четырёх наборах и шести моделях
MAE снижается в среднем на 14{,}5\% относительно лучшего сравнения, а в
offline-replay политика с прогнозом тратит на 21{,}3\% меньше токенов, чем
фиксированный бюджет. Важное наблюдение: разброс расхода между прогонами
превышает порядок величины, поэтому сравнивать факт с точечной оценкой
бессмысленно.

\textbf{Бюджетная осведомлённость — это не одно умение.}
Обзор-эссе~\cite{mahendrakar2026three} разделяет четыре решения:
соблюдение чужого бюджета, наблюдаемость остатка, прогноз собственного
расхода и распределение известного бюджета. Соблюдение в тестируемых
диапазонах (512--3600 токенов) почти решено. Наблюдаемость — нет: ReAct-
агент делает 13{,}77 поисковых вызова при бюджете 30 и 14{,}24 при бюджете
100; плагин, который просто держит остаток в контексте, даёт 12{,}8\% на
BrowseComp при бюджете 10 против 12{,}6\% у немодифицированного агента при
бюджете 100, тратя на 40{,}4\% меньше вызовов и на 31{,}3\% меньше денег.
Главный критический вывод: опубликованные выигрыши от «умного
распределения» часто измерены по сигналу, который стал известен после
факта, а разброс между прогонами вообще не входит в оценку. Это прямой
аргумент, что коррекция должна опираться на причинный, доступный заранее
сигнал, а не на посленаблюдение.

\textbf{Бюджетная осведомлённость агента плохо калибрована.}
BAGEN~\cite{lin2026bagen} формализует её как прогрессивную интервальную
оценку: на каждом шаге агент предсказывает верхнюю и нижнюю границы
остатка и должен предупредить, когда завершение маловероятно. На четырёх
средах и пяти агентах корреляция между способностью решать и бюджетной
осведомлённостью всего $r=0{,}35$; frontier-модели оптимистичны и
продолжают тратить на задачах, которые уже вряд ли решатся. Ранняя
остановка экономит 28--64\% токенов на провальных траекториях, но покрытие
интервалов после SFT+RL не превышает 47\% — калибровка остаётся трудной.

\textbf{Тяжёлые хвосты требуют резервирования, а не среднего.}
CRO~\cite{shokrnezhad2026cro} формулирует задачу оркестратора: зарезервировать
вычисления и полосу так, чтобы покрыть фактический спрос с заданной
вероятностью. Потребление агентского пайплайна неизвестно в момент запуска
и имеет тяжёлые многомодальные распределения. Conformal-резервирование по
скользящему окну прошлых трасс даёт конечновыборочное покрытие и примерно
вдвое снижает перерезервирование по сравнению с границей Чебышёва.

Вывод группы: расхождение по стоимости обнаруживается лучше всего, если
есть распределительный прогноз и обновление по ходу. Точечная оценка и
среднее по прогонам для этой задачи не годятся.

\subsection{Внутренние сигналы: уверенность, неопределённость, аномалии}

Вторая группа ищет сигнал не в стоимости, а в самом ходе решения: уверенность
модели, согласованность траектории, отклонение от нормального поведения.

\textbf{Агенты переоценивают успех, и это измеримо.}
Kaddour и соавторы~\cite{kaddour2026overconfidence} опрашивают агентов до,
во время и после исполнения. Агенты, которые успешны в 22\% случаев,
предсказывают себе 77\% успеха. Неожиданный результат: оценка \emph{до}
исполнения, при строго меньшей информации, различает успех и провал лучше
стандартного посленаблюдения. Лучше всего калибрует адверсариальная постановка
«ищи баги». Отсюда практическое следствие: спрашивать модель «получится
ли» после прогона — не лучший способ; полезнее спрашивать до и в
контр-адверсариальной формулировке.

\textbf{Калибровку можно улучшить процедурой, а не только обучением.}
HTC~\cite{zhang2026htc} извлекает макро- и микродинамику всей траектории и
калибрует уверенность интерпретируемой моделью; на восьми бенчмарках и
разных фреймворках он лучше baseline по ECE, а на OOD GAIA даёт лучшую
калибровку без дообучения. Confidence
Dichotomy~\cite{xuan2026dichotomy} показывает, что тип инструмента меняет
калибровку: evidence-инструменты (поиск) вызывают переоценку из-за шума,
verification-инструменты (интерпретатор) заземляют рассуждение; RL-
дообучение на совместную точность и калибровку переносится на веб и
математику.

\textbf{Неопределённость теряется на передаче.}
Позиционная работа~\cite{shi2026laundering} описывает «отмывку
уверенности»: хрупкое промежуточное решение передаётся следующему
компоненту как чистый артефакт, и локальная неопределённость превращается
в системную ошибку. Это важно для коррекции: если оценка ходит между
планировщиком, исполнителем и оценщиком, она должна нести свою
неопределённость, иначе коррекция будет опаздывать.

\textbf{Провал предсказуем по префиксу.}
PrefixGuard~\cite{huang2026prefixguard} строит обучаемые события и
предсказывает риск провала на горизонте; на 24 комбинациях
бенчмарк--горизонт AUPRC выше сильнейшего baseline на 12{,}0--21{,}9 п.п.,
recall при 10\% FAR — 56{,}1--93{,}0\%. Детерминированный автомат сжимает
историю в рискованные пути; на одной замороженной модели проверка пяти
спецификаций находит нарушения в 21 из 265 успешных прогонов. То есть
формально успешный прогон может нарушать контракт — это отдельный сигнал
ошибки, не видимый по финальному результату.

Doomed from the Start~\cite{ruan2026doomed} идёт дальше: линейные пробы на
скрытых состояниях предсказывают провал с первого раунда взаимодействия,
а каскад с контролем recall даёт гарантию, что успешные эпизоды не будут
отсечены. На TextCraft и WebShop с тремя моделями экономия вычислений
составляет 1{,}5--8{,}8 раза при recall 90\%; токены снижаются на 60{,}2\%
и 54{,}9\% при recall 90\% и на 45{,}0\% и 41{,}5\% при 95\%. Поведенческие
признаки без скрытых состояний работают заметно хуже.

\textbf{Неопределённость накапливается по траектории.}
RUPA~\cite{ma2026rupa} представляет историю как граф рассуждений,
инструментов и обратной связи и распространяет неопределённость по нему;
на $\tau^2$, Terminal-Bench-2 и GAIA с шестью LLM он обнаруживает провал
раньше локальных UQ-методов. Automata from Agent
Traces~\cite{cho2026automata} сворачивает корпус трасс в конечный автомат
на 7--43 состояния; по частичному следу онлайн-монитор ранжирует провальные
прогоны выше успешных, AUROC до 0{,}94, и запускает раннюю остановку.

\textbf{Аномалии и ошибки шага.}
MASC~\cite{shen2025masc} в неразмеченном режиме оценивает аномальность шага
через реконструкцию следующего исполнения и прототипы нормальных шагов;
на Who\&When прирост AUC-ROC до 8{,}47\%, затем агент-корректор исправляет
выход до передачи дальше. TrajAD~\cite{liu2026trajad} формулирует
траекторную аномалию с локализацией для rollback-and-retry и показывает,
что общие LLM плохо локализуют аномалии, а специализированный верификатор
лучше. Agent Vitals~\cite{niedringhaus2026vitals} доводит это до
производственных детекторов loop, stuck, confabulation, thrash и runaway
cost; на 1\,494 трассах все пять проходят порог precision $\ge 0{,}80$ и
recall $\ge 0{,}75$ по нижней границе Уилсона.

Вывод группы: сигналов много, и они информативны. Но почти все работы
измеряют качество классификации (AUROC, AUPRC, precision/recall), а не
пользу решения. Модель, которая хорошо ранжирует провалы, может быть
бесполезна, если цена вмешательства выше выигрыша.

\subsection{Поведенческие симптомы: циклы, дрейф, потеря прогресса}

Третья группа смотрит не на уверенность, а на то, как меняется поведение:
агент зациклился, перестал приближаться к цели, начал «думать» вместо
действия.

\textbf{Overthinking — измеримый симптом неверной оценки.}
Cuadron и соавторы~\cite{cuadron2025overthinking} анализируют 4\,018
траекторий на SWE-bench Verified и выделяют analysis paralysis, rogue
actions и premature disengagement. Чем выше score overthinking, тем ниже
успех; выбор решения с меньшим score даёт почти +30\% успеха и −43\%
стоимости. Это одновременно сигнал (поведение отклонилось) и действие
(выбрать более короткую траекторию). Path deviation
monitoring~\cite{guan2026deviation} использует high-entropy transition
tokens как индекс отклонения от правильного пути и по нему завершает
рассуждение; real-time progress~\cite{raaschou2025progress} учится
предсказывать прогресс 0--100\% с MAE $0{,}161$ на математике.

\textbf{Агент видит бесполезность, но не действует.}
Это, пожалуй, самый важный результат для нашего вопроса. В
ImpasseBench~\cite{chen2026stop} на 19\,371 траектории восьми frontier-
моделей почти каждый прогон, которому нужна внешняя помощь, собирает
достаточные доказательства, но треть всё равно исчерпывает бюджет.
Escalation-инструмент снижает исчерпание с 34{,}0\% до 6{,}7\%, а внешний
монитор почти втрое повышает успех после восстановления среды. В
эксперименте с отказами источника~\cite{zhang2026judged} семь агентов
называют результаты бесполезными в 97--100\% случаев, но редко на этом
останавливаются; правило «остановиться после пяти бесполезных результатов
подряд» помогает, а бюджет в промпте сдвигает остановку 7--8B моделей к
дедлайну. EnoughBench~\cite{almamoori2026enough} формализует это как
normalized sufficiency regret и stop-time calibration: gpt-4o-mini в
reason-act режиме оказывается ниже случайного уровня, и вся
мискалибровка — переоценка; uncertainty-gated hybrid оптимален при
высокой точности сигнала ($\phi \ge 0{,}95$), а при падении точности
начинает собирать лишнее.

\textbf{Точность триггера важнее силы вмешательства.}
В работе о застрявших кодовых агентах~\cite{silva2026trigger} один и тот же
совет без точного триггера даёт −2 п.п. на 200 задачах (7 спасений, 11
регрессий), а точный триггер — пять неудач подряд и сходство stderr не
менее 60\% — убирает 8 из 11 регрессий и меняет итоговый эффект с −4 на +4.
Это прямой аргумент против «вмешиваться при первом сигнале».

Вывод группы: поведенческий симптом — хороший кандидат в сигнал, но сам
по себе он не решает, что делать. Нужна политика с ценой ложной тревоги.

\subsection{Как решить, что оценка неверна}

Между «сигнал сработал» и «оценка неверна» стоит статистическое решение.
Классическая теория даёт для него готовый язык.

\textbf{Последовательное обнаружение изменения.}
Обзор~\cite{xie2021quickest} систематизирует quickest change detection:
наблюдения приходят последовательно, нужно как можно раньше обнаружить
смену распределения при ограничении на ложные тревоги. CUSUM и его
обобщения дают явный компромисс «задержка обнаружения — частота ложных
тревог». Это ровно то, чего не хватает большинству работ по агентам: они
сообщают AUROC, но не говорят, сколько шагов пройдёт до тревоги и сколько
успешных прогонов будет прервано.

\textbf{Гарантия recall вместо порога.}
Doomed~\cite{ruan2026doomed} применяет distribution-free калиброванный
детектор на каждом раннем раунде и оптимизирует бюджеты recall по раундам;
каскад замораживается и сертифицируется на независимых данных, что даёт
пост-селекционную гарантию recall. Это пример правильной постановки:
пользователь задаёт допустимую долю потерянных успешных прогонов, а не
подбирает порог на глаз.

\textbf{Калибровка интервалов — отдельная проблема.}
BAGEN~\cite{lin2026bagen} показывает, что даже после обучения покрытие
интервалов остатка не выше 47\%. CRO~\cite{shokrnezhad2026cro} лечит это
conformal-методом с конечновыборочным покрытием. Отсюда правило: если
прогноз — интервал, проверять надо покрытие, а не MAE.

\textbf{Точность монитора не равна пользе.}
Jain и Vedam~\cite{jain2026monitor} аналитически показывают, что монитор с
фиксированной точностью 90\% может дать от −5 до +35 п.п. к успеху в
зависимости от вероятности успеха заменяющей попытки при остатке бюджета.
На 194 прогонах по 19 задачам SymPy и ретроспективном аудите 130 трасс
первоначальный выигрыш от перезапуска не воспроизвёлся, и ни одно из 103
живых сравнений не прошло сертификацию. Вывод: оценивать монитор надо по
решающей ценности (decision value), а не по accuracy/AUROC.

\textbf{Триггер по противоречию, а не по таймеру.}
BayesBeliefAgent~\cite{goel2026bayes} перепланирует только тогда, когда
действия партнёра прямо противоречат выведенному убеждению; на Overcooked
это даёт на порядок меньше перепланирований, чем эвристики, и сокращает
belief-action gap. Тот же принцип применим к оценке: корректировать не по
каждому отклонению, а когда накопленное свидетельство противоречит
модели.

Синтез: правило обнаружения должно быть последовательным, с явными
параметрами «задержка» и «ложная тревога», и оцениваться по ожидаемой
выгоде действия, а не по площади под ROC-кривой.

Формально минимальное правило выглядит так. Пусть к шагу $t$ есть прогноз
остатка стоимости с медианой $m_t$ и масштабом $s_t$, а наблюдение даёт
остаток $c_t$. Нормированный остаток $r_t = (c_t - m_t)/s_t$ подаётся в
CUSUM:
\begin{equation}
  S_t = \max\bigl(0,\; S_{t-1} + r_t - k\bigr), \qquad
  \text{тревога при } S_t > h .
\end{equation}
Параметры $k$ и $h$ подбираются на исторических трассах так, чтобы доля
ложных тревог не превышала $\alpha$, а задержка обнаружения была
минимальной. После тревоги выбирается действие
$a^\ast = \arg\max_a \mathbb{E}[\text{выгода} \mid \text{свидетельство}]$,
где в множество $a$ входят «продолжать», «сузить», «перепланировать»,
«остановить» и «эскалировать», а выгода считается на остаток бюджета с
учётом стоимости самого вмешательства. Порог $h$ — не константа
предметной области: он калибруется на своих трассах и пересматривается
после каждого прогона.

\subsection{Что делать после обнаружения}

Обнаружение — только половина. Работы этой группы предлагают действия:
обновить прогноз, перепланировать, остановиться, эскалировать,
перераспределить бюджет.

\textbf{Обновить прогноз.}
TokenCast~\cite{ouyang2026tokencast} обновляет кумулятивную стоимость по
новым сегментам; BAGEN~\cite{lin2026bagen} обновляет интервал остатка;
Kalman-подход~\cite{arnold2021kalman} обновляет прогресс процесса по
текущей позиции без журнала (среднее отклонение около 5\%);
байесовский EAC~\cite{caron2012bayesian} даёт интервал итоговой стоимости
с учётом экспертных знаний. Общее: обновлять надо распределение, а не
одно число, и учитывать накопленную ошибку.

\textbf{Перепланировать, но с бюджетом на само перепланирование.}
BRACE~\cite{liu2026brace} формулирует перепланирование как бюджетный контур:
решение «перепланировать или нет», выбор режима, явный токен-бюджет и SLO
по латентности. На Habitat, RoboFactory и AirSim это снижает токены на
replan на 62--92\% и нарушения SLO с 85{,}5--100\% до 4{,}7--50\%, а в
сложной среде даёт 80\% успеха там, где открытый цикл и замороженный план
проваливаются. Should I Replan?~\cite{zahradka2026replan} оценивает
\emph{выгоду} перепланирования по признакам состояния и восстанавливает до
94{,}6\% достижимого снижения задержки. То есть перепланирование — не
бесплатное действие, и его надо сравнивать с продолжением.

\textbf{Восстанавливаться после отказов инструментов.}
ToolMaze~\cite{zhu2026toolmaze} показывает, что неявные семантические
сбои хуже явных: perturbation recovery rate падает примерно на 37\%,
сложные топологии ведут к бесполезным циклам, а отказоустойчивость растёт
с масштабом модели в 3{,}66 раза медленнее, чем базовая задача.
PlanBench-XL~\cite{liu2026planbench} на 327 задачах и 1\,665 инструментах
показывает падение GPT-5.4 с 51{,}90\% до 11{,}36\% при жёсткой
блокировке; труднее всего, когда сбой не даёт явного сигнала. Значит,
детектор должен ловить не только явные ошибки, но и «тихое» расхождение
ожидаемого и фактического результата инструмента.

\textbf{Остановиться или эскалировать.}
Doomed~\cite{ruan2026doomed} превращает прогноз провала в каскад ранней
остановки с гарантией recall. ImpasseBench~\cite{chen2026stop} показывает,
что escalation-инструмент и внешний монитор резко снижают исчерпание
бюджета, но тикет эскалации часто не содержит прогресса и статуса
проверки — то есть «остановиться» и «передать» это разные действия.
EnoughBench~\cite{almamoori2026enough} даёт метрику достаточности
доказательств и штраф за лишний сбор. BAGEN~\cite{lin2026bagen} показывает,
что ранняя остановка экономит 28--64\% токенов на провальных траекториях.

\textbf{Перераспределить бюджет.}
BAVT~\cite{li2026bavt} использует остаток бюджета как показатель степени
при выборе узла поиска и residual value predictor вместо абсолютной
самооценки; при жёстком бюджете он превосходит baseline с четырёхкратным
ресурсом. TRIAGE~\cite{nazi2026triage} показывает, что проспективное
планирование пула задач под общий бюджет — отдельная способность, и
разрыв с oracle велик. CRO~\cite{shokrnezhad2026cro} резервирует ресурсы
заранее по прошлым трассам.

\textbf{Сохранять неопределённость при передаче.}
Если коррекция выполняется несколькими компонентами, хрупкая оценка не
должна превращаться в «чистый» артефакт на handoff~\cite{shi2026laundering}.
Иначе следующий компонент будет over-trust-ить уже исправленную оценку.

\textbf{Два контура коррекции.}
Внутри прогона полезны residual state и журнал зависимостей
(TRACE-FIX~\cite{huang2026tracefix}), который отделяет незакрытые
требования от ярких поверхностных ошибок и резервирует бюджет на ремонт.
Между прогонами работают атрибуция и обучение: PrefixGuard, MASC и TrajAD
дают разметку, по которой можно перекалибровать детектор. Без второго
контура априорная оценка никогда не станет точнее.

\subsection{Пределы самокоррекции}

Соблазнительно поручить коррекцию самой модели: «заметь ошибку и исправь».
Литература говорит, что это ненадёжно.

CorrectBench~\cite{tie2025correctbench} проверяет внутреннюю, внешнюю и
дообученную самокоррекцию на commonsense, математике и коде. Самокоррекция
помогает на сложных задачах, смешение стратегий даёт ещё прирост, но
снижает эффективность; reasoning-модели получают мало, а время растёт;
простой CoT конкурентоспособен. The Illusion of
Insight~\cite{daliberti2026illusion} на 1M+ трасс показывает, что
mid-reasoning сдвиги редки, не учащаются с обучением и редко улучшают
точность; внешние сдвиги при высокой энтропии помогают, но это не
внутренняя самокоррекция. Lost at the Beginning~\cite{liao2025lost}
показывает, что ошибка в первом шаге доминирует, а resampling первых шагов
снижает стоимость до 70\% без потери точности. Synthetic error
injection~\cite{wu2025synthetic} не учит самокоррекции: модель часто
повторяет исходную ошибку, а сдвиг к on-policy ошибкам ломает коррекцию.

Позитивный результат есть, но он требует структуры. Thought-
ICS~\cite{samanta2026thoughtics} разбивает рассуждение на дискретные
семантические шаги, локализует первую ошибку и откатывается к последнему
верному шагу; при oracle-проверке прирост 20--40\%, автономно — лучше
baseline. Метакогнитивный harness~\cite{cao2026metacognitive} разделяет
мониторинг и рассуждение: FOK до решения и JOL после попытки превращаются
в управляющий сигнал; pooled accuracy растёт с 48{,}3 до 56{,}9 без
дообучения. MUSE~\cite{valiente2024muse} учит агента оценивать
компетентность и выбирать стратегию в незнакомой среде.

Вывод: сигнал самооценки полезен, но только как вход в явную процедуру.
Коррекция должна быть внешней, структурной и проверяемой; «попросить
модель подумать ещё раз» — не метод.

\subsection{Классические аналоги}

Задача «заметить, что прогноз разошёлся с фактом, и пересчитать» старше
LLM. Три классические области дают готовые решения.

\textbf{Мониторинг бизнес-процессов.} Обзор~\cite{verenich2019remaining}
сравнивает 16 методов предсказания остаточного времени на 17 реальных
наборах. Kalman-подход~\cite{arnold2021kalman} предсказывает прогресс
запущенного процесса без журнала: 250 оценок, среднее отклонение около
5\%. Логика та же, что нужна харнессу: есть запущенный кейс, есть прогноз
остатка, есть отклонение, по которому принимают решение о ресурсах.

\textbf{Earned value management.} EVM даёт формальный язык для «факт
против плана»: planned value, earned value, actual cost, cost performance
index, estimate at completion. Байесовский EAC~\cite{caron2012bayesian}
строит доверительный интервал итоговой стоимости, duration
EAC~\cite{sackey2020duration} прогнозирует срок в единицах времени, а
анализ~\cite{almeida2021evm} показывает, что детерминированная EVM
наследует ошибку измерения прогресса. Для нас важны два урока: (1)
сравнивать надо не план и факт, а план, факт и неопределённость обоих;
(2) прогресс измеряется с ошибкой, и эта ошибка входит в оценку.

\textbf{Раннее предупреждение о превышении.} Catana и
Florescu~\cite{catana2026resampling} на 4\,329 issue Apache (minority
2{,}2\%) показывают, что почти идеальное качество при полном наборе
признаков — утечка цели, а SMOTE ухудшает cost-sensitive качество; веса
классов надёжнее. Gou~\cite{gou2026earlywarning} на Promise dataset
получает F1 $0{,}876$ и обнаруживает риски в среднем на 2{,}1 итерации
раньше. Урок: раннее предупреждение возможно, но качество сильно зависит
от того, нет ли в признаках будущей информации.

\textbf{Почему обнаруженное не корректируют.} Galorath~\cite{galorath2015estimates}
объясняет систематический оптимизм и стратегическое искажение оценок;
спасает outside view — исторические метрики и параметрические модели.
Smith и Keil~\cite{smith1999badnews} описывают, как плохие новости не
доходят до решающих, а эскалация обязательств усиливает провал; как
минимум четверть проектов неуспешна. Это не статистика, а организационный
механизм: даже хороший детектор не поможет, если сигнал не превращается в
решение.

\subsection{Что следует из сравнения групп}

Три слоя видны отчётливо.

\begin{enumerate}
  \item \textbf{Сигнал.} Здесь больше всего работ и самые сильные
        результаты: провал предсказуем по префиксу, скрытым состояниям,
        графу траектории, автомату, аномалиям шага. AUROC/AUPRC высокие,
        задержка измеряется шагами.
  \item \textbf{Решение.} Здесь работ меньше, и они показывают, что
        accuracy недостаточна: нужна решающая ценность, учёт цены
        вмешательства и вероятности успеха замены. Точные триггеры
        выигрывают у частых.
  \item \textbf{Действие.} Здесь всё ещё фрагментарно: перепланирование,
        остановка и перераспределение изучаются отдельно, редко с общим
        бюджетом и почти никогда с обратной связью на оценщик.
\end{enumerate}

Именно разрыв между слоями и есть пробел.

\section{Сравнение}

Таблица~\ref{tab:sources} сводит включённые работы. Строки сгруппированы
по типу вклада: сигнал, решение, действие, самокоррекция, классический
аналог.

{\scriptsize
\setlength{\tabcolsep}{3pt}
\renewcommand{\arraystretch}{1.12}
\begin{longtable}{@{}>{\raggedright\arraybackslash}p{0.65cm}>{\raggedright\arraybackslash}p{2.9cm}>{\raggedright\arraybackslash}p{2.15cm}>{\raggedright\arraybackslash}p{3.0cm}>{\raggedright\arraybackslash}p{3.6cm}>{\raggedright\arraybackslash}p{2.4cm}@{}}
\caption{Включённые работы.}\label{tab:sources}\\
\toprule
ID & Работа & Что даёт & Метод & Данные и результат & Ограничение \\
\midrule
\endfirsthead
\toprule
ID & Работа & Что даёт & Метод & Данные и результат & Ограничение \\
\midrule
\endhead
\bottomrule
\endfoot

S01 & Lin et al.~\cite{lin2026bagen} & сигнал и действие & интервалы остатка бюджета & 4 среды, 5 агентов; ранняя остановка −28--64\% токенов; покрытие $\le 47\%$ & калибровка слабая \\
S02 & Ouyang et al.~\cite{ouyang2026tokencast} & прогноз стоимости & композиция сегментов & 96 комбинаций; MAE $-14{,}5\%$; $-21{,}3\%$ токенов & не выбирает действие \\
S03 & Kaddour et al.~\cite{kaddour2026overconfidence} & сигнал & опрос до/во время/после & 22\% успеха против 77\% прогноза; pre лучше post & вербальные оценки \\
S04 & Zhang et al.~\cite{zhang2026htc} & сигнал & process-калибровка & 8 бенчмарков; лучший ECE на GAIA & не стоимость \\
S05 & Xuan et al.~\cite{xuan2026dichotomy} & сигнал & RL-калибровка & evidence vs verification; переносимость & tool-use \\
S06 & Shi et al.~\cite{shi2026laundering} & передача & латентная неопределённость & описана «отмывка» & позиция \\
S07 & Huang et al.~\cite{huang2026prefixguard} & сигнал & GRU + автомат & AUPRC $+12{,}0$--$21{,}9$ п.п.; recall 56--93\% & нужны события \\
S08 & Ruan et al.~\cite{ruan2026doomed} & сигнал и остановка & пробы + каскад recall & экономия в 1{,}5--8{,}8 раза; $-60{,}2\%$ токенов & скрытые состояния \\
S09 & Ma et al.~\cite{ma2026rupa} & сигнал & граф неопределённости & $\tau^2$, Terminal-Bench-2, GAIA & нужен граф \\
S10 & Cho et al.~\cite{cho2026automata} & сигнал & FSM по трассам & 12 датасетов; AUROC 0{,}94 & зависит от харнесса \\
S11 & Shen et al.~\cite{shen2025masc} & сигнал и коррекция & реконструкция шага & Who\&When; $+8{,}47\%$ AUC-ROC & MAS \\
S12 & Liu et al.~\cite{liu2026trajad} & локализация & верификатор & общие LLM плохо локализуют & синтетика \\
S13 & Niedringhaus~\cite{niedringhaus2026vitals} & сигнал & 5 детекторов & 1\,494 трассы; precision $\ge0{,}80$ & отчёт Zenodo \\
S14 & Jain, Vedam~\cite{jain2026monitor} & решение & decision value & 90\% точности даёт $-5$--$+35$ п.п.; 0/103 & только restart \\
S15 & Silva Gasca~\cite{silva2026trigger} & решение & точный триггер & бессистемно $-2$ п.п.; точный триггер $-4\to+4$ & малые $n$ \\
S16 & Chen et al.~\cite{chen2026stop} & решение и действие & ImpasseBench & 19\,371 трасса; 34{,}0$\to$6{,}7\% & инъекция блокеров \\
S17 & Zhang et al.~\cite{zhang2026judged} & разрыв & правило остановки & 97--100\% суждений, редко остановка & поиск \\
S18 & Almamoori~\cite{almamoori2026enough} & решение & sufficiency regret & 4 трека; переоценка & бенчмарк \\
S19 & Liu et al.~\cite{liu2026brace} & действие & бюджетный replan & $-62$--$92\%$ токенов; SLO 4{,}7--50\% & embodied \\
S20 & Zahrádka et al.~\cite{zahradka2026replan} & действие & регрессор выгоды & 12\,000 экспериментов; 94{,}6\% & MAPF \\
S21 & Zhu et al.~\cite{zhu2026toolmaze} & действие & бенчмарк отказов & PRR $-37\%$; масштаб помогает в 3{,}66 раза медленнее & нет политики \\
S22 & Liu et al.~\cite{liu2026planbench} & действие & блокировки инструментов & 51{,}90$\to$11{,}36\% & retail \\
S23 & Goel et al.~\cite{goel2026bayes} & действие & байесовский триггер & на порядок меньше replan & координация \\
S24 & Wang et al.~\cite{wang2026epc} & действие & epistemic calibration & $+9{,}75\%$ успеха & без токенов \\
S25 & Huang, Lim~\cite{huang2026tracefix} & действие & residual state & лучше локализация и ремонт & нет чисел \\
S26 & Valiente, Pilly~\cite{valiente2024muse} & действие & метакогниция & лучше OOD, чем RL/prompt & не бюджет \\
S27 & Cao et al.~\cite{cao2026metacognitive} & решение & FOK/JOL harness & 48{,}3$\to$56{,}9 & без стоимости \\
S28 & Nazi, Dipta~\cite{nazi2026triage} & решение & проспективный план & большой разрыв с oracle & способность \\
S29 & Mahendrakar~\cite{mahendrakar2026three} & анализ & 4 решения о бюджете & observability даёт 12{,}8\% при бюджете 10 & эссе \\
S30 & Shokrnezhad, Taleb~\cite{shokrnezhad2026cro} & резерв & conformal & вдвое меньше перерезервирования & инфраструктура \\
S31 & Li et al.~\cite{li2026bavt} & действие & budget-conditioned search & при жёстком бюджете лучше 4$\times$ & QA \\
S32 & Ding et al.~\cite{ding2026cta} & действие & cost-uncertainty & лучше стратегии & короткие задачи \\
S33 & Tie et al.~\cite{tie2025correctbench} & самокоррекция & CorrectBench & помогает, но дорого & бенчмарк \\
S34 & d'Aliberti, Horta Ribeiro~\cite{daliberti2026illusion} & самокоррекция & 1M+ трасс & сдвиги редки и слабы & reasoning \\
S35 & Liao et al.~\cite{liao2025lost} & самокоррекция & resampling первого шага & $-70\%$ стоимости & reasoning \\
S36 & Wu et al.~\cite{wu2025synthetic} & самокоррекция & синтетические ошибки & не работает & negative \\
S37 & Samanta et al.~\cite{samanta2026thoughtics} & самокоррекция & Thought-ICS & $20$--$40\%$ при oracle & нужен oracle \\
S38 & Cuadron et al.~\cite{cuadron2025overthinking} & симптом & overthinking score & 4\,018 трасс; $+30\%$ и $-43\%$ & нужна разметка \\
S39 & Guan et al.~\cite{guan2026deviation} & симптом & path deviation & лучше early-exit & reasoning \\
S40 & Raaschou-Jensen et al.~\cite{raaschou2025progress} & прогресс & проба + дообучение & MAE 0{,}161 & математика \\
S41 & Xie et al.~\cite{xie2021quickest} & решение & quickest detection & CUSUM, tradeoff & классика \\
S42 & Verenich et al.~\cite{verenich2019remaining} & аналог & 16 методов & 17 наборов & процессы \\
S43 & Arnold et al.~\cite{arnold2021kalman} & аналог & Kalman & 250 оценок; 5\% & процессы \\
S44 & Caron et al.~\cite{caron2012bayesian} & аналог & байесовский EAC & интервалы EAC & один кейс \\
S45 & Sackey et al.~\cite{sackey2020duration} & аналог & DEAC & точнее EVM по срокам & строительство \\
S46 & Almeida et al.~\cite{almeida2021evm} & аналог & анализ ошибки EV & прогресс измеряется с ошибкой & аналитика \\
S47 & Catana, Florescu~\cite{catana2026resampling} & аналог & resampling & 4\,329 issue; утечка цели & issue-level \\
S48 & Gou~\cite{gou2026earlywarning} & аналог & GA + FDT & F1 0{,}876; на 2{,}1 итерации раньше & датасет \\
S49 & Galorath~\cite{galorath2015estimates} & организация & outside view & оптимизм и стратегическое искажение & не эмпирика \\
S50 & Smith, Keil~\cite{smith1999badnews} & организация & модель эскалации & четверть проектов неуспешна & не количественно \\
\end{longtable}
}

\section{Пробел и вывод}

\subsection{Что не покрыто}

\begin{enumerate}
  \item \textbf{Совместная схема «обнаружить — решить — исправить».} Работы
        дают либо сигнал, либо правило перепланирования, либо обновление
        прогноза. Нет работы, которая для кодового агента при общем бюджете
        токенов соединяет априорную оценку, онлайн-обнаружение расхождения
        с калиброванной ложной тревогой и выбор действия по ожидаемой
        выгоде.
  \item \textbf{Метрика пользы, а не качества классификации.} Большинство
        детекторов сообщает AUROC/AUPRC/accuracy. Лишь единицы измеряют
        сэкономленные токены, прирост успеха или решающую ценность
        (S08, S14, S16, S19). Без такой метрики нельзя выбрать порог.
  \item \textbf{Разброс между прогонами и тяжёлые хвосты.} TokenCast и
        Three Things показывают, что расход одной задачи меняется на
        порядок, а прогнозы часто сравнивают с одним прогоном. Работ,
        которые моделируют распределение и проверяют покрытие, мало
        (S01, S02, S30).
  \item \textbf{Два контура коррекции не замкнуты.} Внутрипрогонное
        перепланирование и межпрогонное обновление оценщика изучаются
        отдельно. Атрибуция ошибок (S07, S11, S12) не превращается в
        перекалибровку априорной оценки.
  \item \textbf{Разные причины смешаны.} Ошибка оценки, ошибка исполнения
        и невыполнимость требуют разных действий, но в литературе они
        часто сводятся к «провалу». Из-за этого детектор нельзя правильно
        обучить: метка «плохо» не говорит, что исправлять.
  \item \textbf{Организационный разрыв.} Даже при наличии сигнала агент
        или harness может не действовать (S16, S17), а вмешательство без
        точного триггера ухудшает результат (S15, S50).
\end{enumerate}

\subsection{Как это влияет на наше решение}

Литература не даёт готового модуля, но даёт последовательность решений.

\begin{enumerate}
  \item \textbf{Хранить прогноз как распределение.} Для каждой задачи —
        квантили токенов, число шагов, вероятность успеха и вероятность
        выполнимости. Точечная оценка не позволяет отличить шум от
        ошибки.
  \item \textbf{Считать остатки, а не абсолютные значения.} На каждом шаге
        сравнивать фактический расход и прогресс с прогнозом; нормировать
        на неопределённость прогноза. Порог по абсолютному остатку бюджета
        (как в BAGEN) — запасной вариант, а не основной сигнал.
  \item \textbf{Применять последовательный тест.} CUSUM или e-process по
        нормированным остаткам с явными параметрами «задержка» и «ложная
        тревога». Требовать накопления свидетельства, а не реагировать на
        один шаг. Проверять покрытие интервалов, а не только MAE.
  \item \textbf{Выбирать действие по решающей ценности.} Сравнивать
        «продолжать», «сузить область», «сменить стратегию»,
        «перепланировать», «остановить», «эскалировать» по ожидаемой
        выгоде на остаток бюджета. Учитывать вероятность успеха
        заменяющей попытки — иначе перезапуск может навредить
        (S14).
  \item \textbf{Ограничить частоту вмешательств.} Требовать точного
        триггера (несколько неудач, сходство ошибок), иначе выигрыш
        сменяется регрессией (S15). Оценивать не только спасённые задачи,
        но и испорченные.
  \item \textbf{Разделить причины.} Вести отдельные сигналы для ошибки
        оценки (расход/прогресс), ошибки исполнения (аномалия шага) и
        невыполнимости (доказательства блокировки). Действия для них
        разные.
  \item \textbf{Замкнуть два контура.} Внутри прогона — обновление и
        действие; после прогона — атрибуция и перекалибровка оценщика.
        Логировать сигнал, действие, исход и стоимость вмешательства.
  \item \textbf{Не полагаться на вербальную самокоррекцию.} Использовать
        её как один из признаков, но решение принимать по внешнему
        сигналу и явному правилу (S33--S37).
  \item \textbf{Добавить outside view.} Для холодного старта — исторические
        метрики похожих задач и ссылочный класс, а не только оценка
        текущей траектории (S49). Это корректирует систематическое
        смещение, которое онлайн-сигнал может не успеть заметить.
  \item \textbf{Оценивать стоимость мониторинга.} Сигнал должен стоить
        доли процента бюджета. Скрытые состояния и повторные прогоны
        дороги; статические признаки, остатки и дешёвые пробы — приемлемы.
\end{enumerate}

Проверять это надо экспериментом: собрать траектории на доступных задачах,
разметить момент, когда оценка стала заведомо неверной, и сравнить три
политики — «не корректировать», «порог по остатку бюджета» и
«последовательный тест + выбор действия». Метрики: задержка обнаружения,
доля ложных тревог, калибровка обновлённого прогноза, сэкономленные токены,
прирост успеха при фиксированном бюджете. Запись эксперимента создаётся в
\texttt{research/experiments/} до запуска. Если прирост устойчив, писать
спецификацию изменения \texttt{04\_budget\_and\_selection.md}.

\subsection{Короткий ответ}

Оценка неверна не тогда, когда факт отличается от числа, а тогда, когда
накопленное расхождение выходит за пределы неопределённости прогноза и
меняет оптимальное действие. Наблюдать надо за остатком стоимости,
прогрессом к цели, доказательствами выполнимости и поведенческими
симптомами (циклы, overthinking, повторяющиеся отказы). Решение принимать
последовательным тестом с явным компромиссом «задержка — ложная тревога»,
а не по одному шагу. Исправлять — на трёх уровнях: обновить распределение
прогноза, сменить действие (продолжать, сузить, перепланировать,
остановить, эскалировать), перекалибровать оценщик по итогам прогона.
Ключевая метрика — решающая ценность, а не точность монитора. Полагаться
на самокоррекцию модели нельзя; нужен внешний сигнал, структура и явное
правило. Готовой интегрированной схемы для кодового агента при общем
бюджете литература не даёт — это и есть пробел, который надо закрывать
экспериментом.

\section{Ограничения обзора}

\begin{itemize}
  \item Две базы вместо четырёх: Semantic Scholar и Google Scholar вернули
        HTTP 403. Возможны пропуски, особенно в индустриальных отчётах.
  \item Просматривались верхние 1--30 записей на запрос, а не вся выдача.
  \item Для большинства работ данные извлечены из аннотаций; полные тексты
        читались выборочно. Числа могут отличаться от итоговых таблиц.
  \item Препринты 2026 года не рецензированы; часть работ по агентному
        мониторингу может измениться.
  \item Метрики несопоставимы: AUROC, AUPRC, ECE, MAE, сэкономленные
        токены, доля успеха. Мета-анализ невозможен.
  \item Классические работы по управлению проектами опираются на
        строительные и IT-проекты, а не на агентов; переносится сигнал, а
        не вывод.
  \item Обзор не проверяет воспроизводимость включённых работ и не
        оценивает риск публикационного смещения.
\end{itemize}

\bibliographystyle{plain}
\bibliography{refs}

\end{document}
